\ConceptDivider{01}{Posterior predictive decisions}{A probability becomes useful only when it is paired with an action and a cost.}

\begin{frame}{Probability is an input to a decision, not the decision itself}
  \small
  \begin{BayesTable}{>{\raggedright\arraybackslash}p{.23\linewidth} X X}
    \BayesTableHeader \BayesTableHead{Question} & \BayesTableHead{Posterior answers} & \BayesTableHead{Policy answers} \\
    What may happen? & A distribution over outcomes. & Nothing yet. \\
    What should we do? & Expected outcomes under each action. & Which trade-off is acceptable. \\
    Who owns the choice? & The model author explains uncertainty. & The decision owner sets the cost and constraint. \\
  \end{BayesTable}
  \vspace{.24cm}
  \begin{KeyTakeaway}{Retain uncertainty until the action boundary}
    Collapsing a posterior too early hides the sensitivity of the final decision to costs, capacity and risk tolerance.
  \end{KeyTakeaway}
\end{frame}

\begin{frame}{Close the learning loop}
  \begin{ExerciseCard}{One-minute retrieval prompt}
    Describe the difference between a prior and a policy in one sentence. Then name one assumption behind the posterior you would want to test first.
  \end{ExerciseCard}
\end{frame}
